\section{Introduction}

Ce document constitue une analyse fonctionnelle des modes de marche et
d'arrêt de la cellule robotique support du module \emph{Automatisme pour la
Robotique}. Il complète le TD de conception (dossier \texttt{01\_conception})
auquel on renverra pour la présentation détaillée du système, de ses
constituants et de ses contraintes ; seul un très bref rappel en est fait
ici.

\begin{UPSTIinfor}{Rappel du système}
La cellule associe un robot SCARA, un convoyeur à palettes régulé par des
bloqueurs, un poste \textbf{guichet} (interface entre le convoyeur et le
robot, de \textbf{capacité unitaire}), une IHM et un automate Schneider
M340. Le besoin consiste à faire tremper des gobelets, transportés par des
palettes, dans l'une des quatre cuves de traitement disponibles. Pour une
présentation complète, se reporter à \texttt{01\_conception/sources/\allowbreak
presentation\_systeme.tex}.
\end{UPSTIinfor}

\subsection{Objet du document}

Ce document est un document de référence à consulter pour connaitre les attendus 
de fonctionnement de la cellule robotisée.

Pour chaque mode retenu, ce document précise :
\begin{itemize}
    \item les \textbf{conditions d'entrée} dans le mode (l'évènement ou la
    condition qui déclenche son activation) ;
    \item le \textbf{comportement attendu} du robot, du convoyeur/guichet
    et de l'IHM durant ce mode (le résultat visé, jamais le mécanisme mis
    en \oe uvre pour l'obtenir) ;
    \item les \textbf{conditions de sortie} du mode, c'est-à-dire vers
    quel(s) autre(s) mode(s) l'installation peut évoluer, et sous quelle
    condition.
\end{itemize}

\subsection{Structure GEMMA retenue}

L'analyse du GEMMA propre à cette installation fait apparaître, pour
la partie de la cellule effectivement étudiée dans ce module, les modes
suivants, répartis selon les trois familles habituelles du GEMMA :

\begin{itemize}
    \item \textbf{Famille D -- Procédures en Défaillance de la Partie
    Opérative}
    \begin{itemize}
        \item Mode \textbf{D1} -- Arrêt d'urgence
        \item Mode \textbf{D2} -- Diagnostic / traitement de la défaillance
    \end{itemize}
    \item \textbf{Famille A -- Procédures d'Arrêt de la Partie Opérative}
    \begin{itemize}
        \item Mode \textbf{A5} -- Préparation pour remise en route après
        défaillance
        \item Mode \textbf{A6} -- Mise en conditions initiales
        \item Mode \textbf{A1} -- Arrêt en conditions initiales
    \end{itemize}
    \item \textbf{Famille F -- Procédures de Fonctionnement}
    \begin{itemize}
        \item Mode \textbf{F1} (couplé à une sortie vers \textbf{A2}) --
        Production normale
        \item Mode \textbf{F4} -- Marche de vérification dans le désordre
        (mode manuel)
    \end{itemize}
\end{itemize}

\subsection{Enchaînement général des modes}

L'automate étant considéré comme repassant systématiquement en mode
\textbf{D1} dès qu'une défaillance de communication avec l'un des
équipements de la cellule (baie robot, contrôleur de l'IHM, contrôleur du convoyeur) est
détectée, l'enchaînement de référence entre les modes retenus est le
suivant :

\begin{center}
\begin{tabular}{cl}
D1 & $\longrightarrow$ D2 (dès que la communication est rétablie) \\
D2 & $\longrightarrow$ A5 (une fois la configuration de l'installation
    assainie) \\
A5 & $\longrightarrow$ A6 (une fois la remise en route autorisée par
    l'opérateur) \\
A6 & $\longrightarrow$ A1 (une fois les conditions initiales atteintes) \\
A1 & $\longrightarrow$ F1 (sur demande de marche automatique) \\
    & $\longrightarrow$ F4 (sur demande de marche manuelle) \\
F1 & $\longleftrightarrow$ A2 (arrêt/reprise en fin de transfert de
    gobelet) \\
F4 & $\longrightarrow$ A6 (sur demande de retour aux conditions
    initiales) \\
\text{tout mode} & $\longrightarrow$ D1 (dès qu'une défaillance de
    communication est détectée) \\
\end{tabular}
\end{center}

Chacun de ces modes est détaillé dans les sections suivantes, dans cet
ordre.
